\documentclass[a4paper,article,oneside,10pt]{memoir}

\usepackage[utf8]{inputenc}
\usepackage[T1]{fontenc}
\usepackage{babel}
\usepackage{kpfonts,microtype}

\usepackage{leftindex,listings,showexpl}

\usepackage[hidelinks]{hyperref}

\frenchspacing

\newcommand\entrydescription[1]{%
	\textlangle\textrm{\textit{#1}}\textrangle
}

\catcode`<=\active

\def<#1>{\entrydescription{#1}}

\lstset{
	language=[LaTeX]TeX,
	basicstyle=\ttfamily\small,
	commentstyle=\itshape\ttfamily\small,
	escapechar=@,
	alsoletter=\\,
	breaklines=true,
	breakindent={0pt},
	captionpos=t,
	pos=r,
	tabsize=2,
	inputencoding=utf8,
	extendedchars=true,
	explpreset={numbers=none,},
}

\newcommand\pkgleftindex{\texttt{leftindex}}

\title{\pkgleftindex\ -- left indices with better spacing (v0.2$\beta$)}
\date{\today}
\author{Sebastian \O rsted (\href{mailto:sorsted@gmail.com}{sorsted@gmail.com})}

\hypersetup{
	pdfauthor={Sebastian \O rsted},
	pdftitle={leftindex -- left indices with better spacing},
}

\begin{document}

\maketitle

\noindent
The package \pkgleftindex\ provides commands for typesetting left indices.
Unlike any other similar package I know of, \pkgleftindex\ also indents the left superscript, producing much better spacing in general:
\[
	\leftindex^a_b{f}
	\qquad\text{(compare to $\leftindex[]^a_b{f}$)}
	.
\]
As \TeX\ provides no metrics for the indentation of left superscripts,
the only way to accomplish this is to measure the negative indentation of the right subscript and then use that value.
This works well for most symbols, with a few exceptions, which can be dealt with by
supplying a few extra arguments.
The package is a generalization of egreg's excellent code found at
\begin{center}
	\url{https://tex.stackexchange.com/questions/557655/low-level-command-for-controlling-left-superscript-spacing/}
\end{center}

\chapter*{The commands}

The main user-level command of~\pkgleftindex\ is called~\lstinline!\leftindex! and has
the syntax
\begin{lstlisting}
\leftindex^{@<left superscript>@}_{@<left subscript>@}{@<symbol>@}
\end{lstlisting}
Both the arguments <left superscript> and <left subscript>
are optional. Let us see it in action:
\begin{LTXexample}
$ \leftindex^a_b {f} $,
$ \leftindex^t {X} $
$ \leftindex_b {J} $,
$ \leftindex_v^u {\Pi} $,
$ \leftindex{A} $,
\[
	\leftindex^1_0 { \int }
		f(x)\, dx
\]
\end{LTXexample}
For the indentaiton of the left superscript, the package by default
uses the negative indentation of the right subscript. This yields good results
most of the time, as seen above. However, there are cases where this goes wrong,
the most notable example being the Greek letter~$ \Gamma $ in many fonts (not in Computer Modern, though). Less serious examples arise with letters like $P$ and $L$:
\begin{LTXexample}
$ \Gamma^a_b $
vs.
$ \leftindex^a_b {\Gamma} $,
$ L^a_b $
vs.
$ \leftindex^a_b {L} $,
$ P^a_b $
vs.
$ \leftindex^a_b {P} $
\end{LTXexample}
To solve this issue, \lstinline!\leftindex! takes two
additional, optional arguments:
\begin{lstlisting}
\leftindex[@<slanting phantom>@][@<height phantom>@]
			^{@<left superscript>@}_{@<left subscript>@}{@<symbol>@}
\end{lstlisting}
The package will then use the <slanting phantom> instead of the <symbol> to calculate the indentation of the superscript. The argument can be left empty, in which
case the left superscript will not be indented at all, which is the right solution for the symbol~$ \Gamma $:
\begin{LTXexample}
$\leftindex[]^a_b {\Gamma}$,
$\leftindex[I]^a_b {L}$,
$\leftindex[I]^a_b {P}$
\end{LTXexample}
By default, the <slanting phantom> is also used for calculating the height of the indices.
If this is not the behaviour you want, you can supply an additional <height phantom>:
\begin{LTXexample}
$\leftindex[][\Bigg|]^a_b{x}$
\end{LTXexample}
I doubt anyone will ever really need this much in practice, but the option is provided for completeness.

Sometimes, having to include the symbol as a mandatory argument can be a burden. For these cases, I also provide the command~\lstinline!\LeftIndex! with the syntax:
\begin{lstlisting}
\LeftIndex[@<slanting phantom>@][@<height phantom>@]
			^{@<left superscript>@}_{@<left subscript>@}
\end{lstlisting}
Here, all arguments are optional.
With this, we can do stuff like:
\begin{LTXexample}
$\LeftIndex^a_b \Gamma $,
$\LeftIndex[]^a_b L $,
$\LeftIndex[I]^a_b$,
$\LeftIndex[I]^a_b$
\end{LTXexample}


I also provide the more primitive, underlying command with four mandatory arguments:
\begin{lstlisting}
\manualleftindex{@<height phantom>@}{@<slanting phantom>@}
			{@<left superscript>@}{@<left subscript>@}
\end{lstlisting}
This one is mainly intended for use in commands and by other packages.
Note that the arguments <height phantom> and <slanting phantom> have
switched places compared to the commands \lstinline!\leftindex! and~\lstinline!\LeftIndex!.
This is because <height phantom> is regarded as a ``more important''
argument than <slanting phantom> from an implementation point of view,
but the user is probably less likely to change it. Therefore,
the user-level commands \lstinline!\leftindex! and~\lstinline!\LeftIndex! have <slanting phantom>
before <height phantom>.
Note also that \lstinline!\manualleftindex! (like \lstinline!\LeftIndex!) does not take a <symbol>;
you will have to write this yourself afterwards:
\begin{LTXexample}
$\manualleftindex{P}{I}{a}{b}$
vs.
$\manualleftindex{P}{I}{a}{b} P$
\end{LTXexample}

\end{document}